\section{Discussion}

\label{sec:discussion}

\subsection{When to Use Code vs. Visual}

Hanzo Flow is most effective for:
\begin{itemize}[leftmargin=1.1em]
    \item Workflows with clear input/output contracts at each step.
    \item Multi-model or multi-agent systems requiring visibility.
    \item Rapid prototyping and iteration.
    \item Teams with mixed technical skill levels.
\end{itemize}

Code-based approaches remain preferable for:
\begin{itemize}[leftmargin=1.1em]
    \item Highly dynamic workflows where structure changes per-input.
    \item Integration with complex existing codebases.
    \item Performance-critical inner loops.
\end{itemize}

\subsection{Limitations}

\begin{enumerate}[leftmargin=1.4em]
    \item \textbf{Expressiveness ceiling}: Some algorithms are awkward to express as dataflow graphs (e.g., recursive algorithms with complex state).
    \item \textbf{Large graphs}: Workflows with 50+ nodes become visually cluttered; subgraph abstraction partially addresses this.
    \item \textbf{Custom nodes}: Extending with custom node types requires TypeScript knowledge.
    \item \textbf{Version control}: Graph-based workflows have less mature diffing/merging than text-based code.
\end{enumerate}

\subsection{Future Work}

\begin{itemize}[leftmargin=1.1em]
    \item \textbf{AI-assisted construction}: Use LLMs to generate workflow graphs from natural language descriptions.
    \item \textbf{Automated optimization}: Apply graph optimization passes (node fusion, reordering, caching) to reduce cost and latency.
    \item \textbf{Collaborative editing}: Real-time multi-user workflow editing with conflict resolution.
    \item \textbf{Marketplace}: Community-contributed workflow templates and custom nodes.
\end{itemize}

